% ===== §6 Symmetry-breaking and the cognition of the unbroken structure =====
\section{Symmetry-Breaking and the Cognition of the Unbroken Structure}
\label{p10:sec:breaking}

\noindent
Two realisations of one structure stand in juxtaposition. The quantum-structural grammar (\xref{p10:sec:grammar}) and the political economy of the reconfigured subject (\xref{p10:sec:political}) are recognisably the same reflexive structure, the reconfiguration of the dynamics by the phenomenon, of the producer by the product, and are irreducibly distinct, neither translatable into the other, each employing a vocabulary the other lacks. This section specifies the status of that situation and addresses the question it generates: given one abstract structure realised in these two distinct forms, the cognition of the unbroken structure, in the absence of the external standpoint excluded by the ontology of \xref{p10:sec:onto-refuses}. The section is the central juncture of the paper, at which the claim to a metalanguage is both most available and most consequential.

\subsection{The disciplines as symmetry-broken phases}
\label{p10:sec:phases}

The relation between the two realisations is specified by symmetry-breaking. There is one abstract, highly symmetric, undifferentiated structure, the reflexive structure in which a generative system's appearance reconfigures its dynamics. Under the conditions of a particular discipline this symmetry breaks, and the structure assumes a determinate, lower-symmetry form: a participatory measurement under the conditions of formal epistemology, the alienation--realisation of the subject under the conditions of political economy, and, as the prior papers established, the dialectic of desire under psychoanalysis, and the principle of reversal-as-the-movement-of-process in classical Chinese cosmology. The operative feature of symmetry-breaking is that, the symmetry once broken, there is no path internal to any broken phase by which the unbroken symmetry is recovered. The symmetric phase is not surveyable from within the asymmetric one; one is always already within a phase, and the unbroken structure is, from within, irrecoverable, accessible only as that which must have broken in order to yield the phase one occupies.

\begin{claim}[The disciplines as symmetry-broken phases of one structure]
Each discipline is not a partial external view of the unbroken structure but a phase that the structure assumes when realised under that discipline's conditions. The absence of a metalanguage, established in the prior paper, obtains in consequence: there is no unbroken phase to occupy, since all discourse proceeds within some discipline's broken vocabulary, from within a determinate phase, the symmetric structure being not a language into which one might revert but precisely that of which every actual language is a breaking. The exclusion of the external standpoint in \xref{p10:sec:onto-refuses} is the ontological statement of this result; symmetry-breaking is its mechanism.
\end{claim}

\subsection{The paradox of the section}
\label{p10:sec:paradox-met}

A paradox is now to be addressed rather than evaded. If one is always already within a broken phase, and the unbroken structure is irrecoverable from within any phase, the possibility of the present section is in question. For the section speaks of the unbroken structure, names it, asserts that the two realisations are phases of it. The objection is that the section thereby claims precisely the external, phase-transcending standpoint it has declared impossible, and that the predicate ``unbroken structure'' constitutes a covert metalanguage, a settling of what has been asserted to be unsettlable.

The paradox is genuine and is not to be dissolved by qualification; it is the substantive difficulty of the undertaking, and the remainder of the section consists in three modes of addressing it without claiming the excluded standpoint. The inadequate response is the claim to a partial overview, that the unbroken structure is dimly apprehended. This retains the standpoint under the appearance of modesty. The adequate response is that the unbroken structure is cognised not by being viewed, from any standpoint however partial, but in three modes, none of which occupies an external position.

\subsection{The unbroken structure as abductively posited}
\label{p10:sec:abduction}

The unbroken structure is not perceived; it is posited, abductively, as the best explanation of a cross-phase correspondence. It is not viewed. What is encountered, from within the phases, is the fact that two disciplines as mutually foreign as quantum measurement and the critique of political economy exhibit one reflexive structure. From that fact the inference proceeds: there is a structure, antecedent to the breaking, of which these are phases, on pain of treating the correspondence as a coincidence of implausible magnitude.

\begin{claim}[The unbroken structure as inference to the best explanation]
The unbroken structure is posited as the inference to the best explanation of the cross-phase isomorphism; it is not given but inferred, and inferred from within the phases. The procedure preserves the constraint the paradox imposes: the unbroken structure is at no point occupied, its features at no point read off directly; one remains within a phase and infers, from the correspondence of the phases, a unity one does not occupy. The positing is defeasible and does not resolve into a view. It is a commitment the phases incur, not a content possessed.
\end{claim}

\subsection{Holonomy as the proposed cross-phase invariant}
\label{p10:sec:invariant}

A more determinate formulation of what is conserved across the phases yields the most consequential, and most carefully delimited, proposal of the paper. Let each discipline be treated as a local gauge, a coordinate chart in which the structure is written in one manner; the passage between disciplinary vocabularies is then a change of chart, a transition function, and the question of the identity of the unbroken structure becomes the question of the invariant under all such changes of chart. In the geometry employed throughout the series there is a distinguished gauge-invariant quantity, one belonging to no single chart yet conserved across all of them, namely the holonomy, the phase accumulated around a loop, on which the good and vicious cycle have turned since the prior paper. The proposal is the following.

\begin{claim}[Holonomy as the cross-phase invariant: a proposal advanced as unsettled]
The distinction between the good and the vicious, the sign of the holonomy, augmentation as against extraction, the spiral as against the circle, is proposed as the invariant conserved across the broken phases: belonging to no single discipline, conserved across all, and capable of serving as the connective among disciplines in virtue of being possessed by none. On this proposal the geometric phase tracked across intimacy, political economy, quantum epistemology, and classical cosmology is not a further discipline but the gauge-invariant core of the unbroken structure, which accounts for the operation of one criterion across all of them. The proposal is advanced as unsettled. To assert it as the established truth of the unbroken structure would be to cash, in a sign, a determination the paper has not established, to commit, at its own central juncture, the settling failure it elsewhere identifies. It is accordingly advanced as the strongest unsettled conjecture of the series, the mode of whose assertion must, on pain of self-contradiction, remain non-possessive.
\end{claim}

\noindent
The restraint is required and is not a weakness of the position. The available move is the assertion that holonomy is the invariant and the central juncture thereby established, with the consequent completion of the series. But the section has established that the unbroken structure is not possessible from any phase; its subsequent possession, under the designation of holonomy, would controvert the section in its conclusion. The conjecture is therefore left standing, at once load-bearing and unsettled, which is the only manner in which a central juncture of this structure can stand. The formal vocabulary of the conjecture, and the conditions on its discharge, are recorded in \xref{p10:app:phase}.

\subsection{The cognition of the unbroken structure as the traversal of the phases}
\label{p10:sec:walk}

The third mode of addressing the paradox is performative rather than asserted. The manner in which the section has cognised the unbroken structure is not by viewing it but by traversing the phases, by passing from the quantum phase to the political-economic phase and registering, in the passage, their correspondence. The cognition of the unbroken structure is the holonomy accumulated by the interpreter who traverses the loop from one phase to another and returns; it is not a proposition delivered concerning the unbroken structure but a phase accumulated in traversing its breakings. This is the poetics of the prior paper applied to the method of the present one: the unbroken structure is not represented, representation being a settling and a claim to the standpoint, but witnessed, by an interpreter who traverses the path through the phases and accumulates, in the traversal, a cognition no summary delivers.

\begin{claim}[Cognition by traversal rather than survey]
The section does not survey the phases from above; it traverses them, and its cognition of their common structure is the phase accumulated in the traversal. The section therefore terminates in no proposition concerning the unbroken structure, but in the traversed recognition that the phases correspond. The form of the section is its content: a refusal to settle the unbroken structure into a thesis, effected in the manner of the section's procedure.
\end{claim}

\subsection{The dialectical-materialist ground of the polyphony}
\label{p10:sec:dialmat}

A misidentification is to be precluded. The insistence on a plurality of frameworks, none sovereign, no metalanguage, a truth obtaining in the dialogue of phases rather than in any one, admits of construal as postmodern relativism, the abandonment of truth for an irreducible plurality of language-games. The position is the contrary, and the distinction is the methodological commitment of the series.

\begin{claim}[The polyphony as dialectical-materialist, not relativist]
The plurality of frameworks is not the relativist thesis that there is no truth but only discourses. It is the dialectical-materialist thesis that truth is real and not yet complete, that it is realised in a historical, material process of generation, such that any single framework is a partial determination, from one historical-disciplinary position, of a totality still in production. No framework possesses the whole, not because there is no whole, but because the whole is unfinished; and the declaration of its capture, in any single framework, is the settling of a process in progress into a fixed doctrine, which is alienation in the register of cognition. The many-framework polyphony of the series is accordingly neither eclecticism nor collage but the method a not-yet-complete totality requires: the available fidelity to a truth still in the course of its own generation. That generative justice must sustain itself through a recursive structure, that the poem refuses to settle, that the dialectic does not close in absolute knowledge, are formulations of one commitment, and it is materialist throughout.
\end{claim}

\noindent
This precludes the construal the form of the series most invites. The objection that the work does not conclude, that it disperses into frameworks rather than synthesising them, misidentifies a method as a deficiency. The refusal to synthesise is not an inability to conclude but a fidelity to a totality that is not concluded, a refusal, on principle, of the settling that would present an unfinished generation as a completed possession. The polyphony is the form truth assumes when truth is still in production. The section has thereby completed its task: it has specified the relation between the disciplines as symmetry-broken phases of one structure; addressed the paradox of cognising the unbroken structure in the absence of an external standpoint, through the abductive positing of the structure, the proposal of holonomy as an unsettled invariant, and the traversal rather than survey of the phases; and established the dialectical-materialist rather than relativist character of the position. The manner in which the unbroken structure is cognised, by traversal rather than survey, by positing rather than possession, by an invariant advanced rather than settled, is the first determination of the practice, being the cognitive form of \emph{wu wei}: a cognition that does not grasp. The practice is the matter of the following section.
